%%%%%%%%%%%%%%%%%%%%%%%%%%%%%%%%%%%%%%%%%
% Beamer Presentation
% LaTeX Template
% Version 1.0 (10/11/12)
%
% This template has been downloaded from:
% http://www.LaTeXTemplates.com
%
% License:
% CC BY-NC-SA 3.0 (http://creativecommons.org/licenses/by-nc-sa/3.0/)
%
% Modified by Jeremie Gillet in November 2015 to make an OIST Skill Pill template
%
%%%%%%%%%%%%%%%%%%%%%%%%%%%%%%%%%%%%%%%%%

%----------------------------------------------------------------------------------------
%	PACKAGES AND THEMES
%----------------------------------------------------------------------------------------

\documentclass{beamer}

\mode<presentation> {

\usetheme{Madrid}

\definecolor{OISTcolor}{rgb}{0.65,0.16,0.16}
\usecolortheme[named=OISTcolor]{structure}

%\setbeamertemplate{footline} % To remove the footer line in all slides uncomment this line
%\setbeamertemplate{footline}[page number] % To replace the footer line in all slides with a simple slide count uncomment this line

\setbeamertemplate{navigation symbols}{} % To remove the navigation symbols from the bottom of all slides uncomment this line
}

\usepackage{graphicx} % Allows including images
\usepackage{booktabs} % Allows the use of \toprule, \midrule and \bottomrule in tables
\usepackage{textpos} % Use for positioning the Skill Pill logo
\usepackage{fancyvrb}
\usepackage{tikz}
\usepackage{hyperref}
\usepackage{listings}

\definecolor{dkgreen}{rgb}{0,0.6,0}
\definecolor{gray}{rgb}{0.5,0.5,0.5}
\definecolor{mauve}{rgb}{0.58,0,0.82}

\lstset{frame=tb,
  language=python,
  aboveskip=3mm,
  belowskip=3mm,
  showstringspaces=false,
  columns=flexible,
  basicstyle={\small\ttfamily},
  numbers=none,
  numberstyle=\tiny\color{gray},
  keywordstyle=\color{blue},
  commentstyle=\color{dkgreen},
  stringstyle=\color{mauve},
  breaklines=true,
  breakatwhitespace=true,
  tabsize=3
}

\newcommand*{\lstitem}[1]{
  \setbox0\hbox{\lstinline{#1}}  
  \item[\usebox0]  
  % \item[\hbox{\lstinline{#1}}]
  \hfill \\
}

%----------------------------------------------------------------------------------------
%	TITLE PAGE
%----------------------------------------------------------------------------------------

\title[Skill Pill]{Skill Pill: Julia} % The short title appears at the bottom of every slide, the full title is only on the title page
\subtitle{Lecture 1: Introduction}

\author{James Schloss \and Valentin Churavy} % Your name
\institute[OIST] % Your institution as it will appear on the bottom of every slide, may be shorthand to save space
{
Okinawa Institute of Science and Technology \\ % Your institution for the title page
\textit{james.schloss@oist.jp}\\
\textit{valentin.churavy@oist.jp} % Your email address
}
\date{July 7, 2017} % Date, can be changed to a custom date

\begin{document}

\setbeamertemplate{background}{\includegraphics[width=\paperwidth]{SPbackground.png}} % Adding the background logo

\begin{frame}
\vspace*{1.4cm}
\titlepage % Print the title page as the first slide
\end{frame}

\setbeamertemplate{background}{} % No background logo after title frame

\addtobeamertemplate{frametitle}{}{% Adding the Skill Pill logo on the title screen after title frame
\begin{textblock*}{100mm}(.92\textwidth,-0.9cm)
\includegraphics[height=0.85cm]{julia.pdf}
\end{textblock*}}

\begin{frame}
  \tableofcontents
\end{frame}
\section{The foreign world, using Julia to reuse prior work}
\begin{frame}[fragile]{Using Fortran and C in Julia}

Julia allows you to use other languages (such as Fortran or C) by using the \texttt{ccall} function:

\begin{lstlisting}
    julia> t = ccall((:clock, "libc"), Int32, ())
    2292761
\end{lstlisting}

Here, we are calling the \texttt{clock} function from the \texttt{libc} library in C.

\end{frame}

\begin{frame}[fragile]{Your legacy code}
Let's say you want to use a simply multiply function in Fortran:
\begin{lstlisting}[language=fortran]
      !! We'll be using subroutines intead of functions 
      subroutine multiply(A, B, C)
          REAL*8 :: A, B, C
          C = A * B
          return
      end
\end{lstlisting}
\vspace{0.5cm}
or C:
\begin{lstlisting}[language=c]
    // Nothing fancy here...
    double multiply(double A, double B){
        return A*B;
    }
\end{lstlisting}
\end{frame}

\begin{frame}[fragile]{Preparing your legacy code}

In order to use your favorite C or Fortran code in Julia, you need to compile it into a library, like so:

\begin{lstlisting}
    gcc -shared -O2 multiply.c -fPIC -o c_multiply.so
    gfortran -shared -O2 multiply.f90 -fPIC -o 
        fortran_multiply.so
\end{lstlisting}

\pause
These will create libraries with all of the necessary functions you could want, but beware:

\begin{center}
\textbf{C and Fortran compilers mangle function names!}
\end{center}
\end{frame}

\begin{frame}[fragile]{Using your legacy code}
There are 3 things to keep. Make sure you 
\begin{enumerate}
\item Have the right mangled name
\item Are using the right type
\item Are using the function correctly.
\end{enumerate}
\pause

\vspace{0.5cm}
For example, in C:
\begin{lstlisting}
    # This function multiplies a and b into c by using the 
    # created C library
    function call_c()
        a = Cdouble(1.0)
        b = Cdouble(3.0)
        c = ccall((:multiply, "/full/path/to/c_multiply"),
            Cdouble,(Cdouble, Cdouble),a,b)
        println(c)
    end
\end{lstlisting}
\end{frame}

\begin{frame}[fragile]{Using your legacy code}
Pointers are okay! For example, in Fortran:

\begin{lstlisting}
    # This function multiplies a and b into c by using 
    # the created FORTRAN library
    function call_fortran()
        a = Cdouble[1.0]
        b = Cdouble[2.0]
        c = Cdouble[0.0]
        ccall((:multiply_, "/full/path/to/fortran_multiply"),
           Void,(Ptr{Float64},Ptr{Float64},Ptr{Float64}),
                 a,b,c)
        println(c[1])
    end
\end{lstlisting}

\pause
More information can be found here: \url{https://docs.julialang.org/en/stable/manual/calling-c-and-fortran-code/}
\end{frame}
\begin{frame}{Support for other languages}
  \begin{description}
    \item[Python] https://github.com/JuliaPy/PyCall.jl
    \item[R] https://github.com/JuliaInterop/RCall.jl
    \item[C++] https://github.com/Keno/Cxx.jl
    \item[Matlab] I have heard rumours of such a thing existing, but the horror $\cdots$
  \end{description}
  \begin{block}{Conclusion}
    Start writing Julia code now without being worried about losing your prior work!
  \end{block}
\end{frame}

\section{The Julia compiler}
\begin{frame}[fragile]{The stages of the compiler}
  \begin{enumerate}
    \item Surface syntax (the code you write)
    \item Desugared AST \lstinline{@code_lowered}
    \item Type-inferred AST \lstinline{@code_typed}
    \item LLVM IR \lstinline{@code_llvm}
    \item Native assembly \lstinline{@code_native}
  \end{enumerate}
\end{frame}

\begin{frame}[fragile,allowframebreaks]{Macros and metaprogramming}
  Julia has powerful metaprogramming facilities that allow you to manipulate the AST (abstract syntax tree) of a Julia program.
  \begin{lstlisting}
  prog = "1 + 1"
  ex1 = parse(prog) # Parse a Julia program and return the expression tree
  # This is equal to manually constructing the expr tree
  ex2 = Expr(:call, :+, 1, 1)
  # Or by using quote
  ex3 = :( 1 + 1) 
  dump(ex1) # Inspect an expr tree
  \end{lstlisting}
  \vspace{-1em}
  \begin{block}{}
    The key point here is that Julia code is internally represented as a data structure that is accessible from the language itself.
  \end{block}
  \vspace{-1em}
  \begin{lstlisting}
  # use eval to execute an expression
  eval(ex1)
  \end{lstlisting}
  \framebreak
  \begin{block}{Macros}
    \lstinline{eval} works at runtime and you should avoid it as much as possible. You can use macros to do syntax transformations statically.
    Macros always start with an \verb|@|.
  \end{block}
  \begin{lstlisting}
  # simplistic implementation of the @assert macro in Julia
  macro assert(ex)
      return :( $ex ? nothing : throw(AssertionError($(string(ex)))) )
  end

  @assert 1 == 1.0
  @assert 1 == 0

  @macroexpand @assert 1 == 0
  \end{lstlisting}
\end{frame}
\section{Performance}
\begin{frame}[fragile]{Methodology}
  \pause
  \begin{block}{Measure first}
    Before you start iterating on your code establish a baseline performance.
    Computers are noisy system so we use the lowest runtime as a metric.
  \end{block}
  \pause
  \begin{itemize}
    \item Check for type-instabilities with \lstinline{@code_warntype}
    \item Measure runtime and allocations with \lstinline{@time}
    \item Benchmark using \lstinline{@btime}, and \lstinline{@benchmark} from \lstinline{BenchmarkTools.jl}
    \item Profiler and \lstinline{ProfileView.jl}
    \item Memory Allocation tracker
  \end{itemize}
  \pause
  Read the performance tips section of the Julia manual \url{https://docs.julialang.org/en/stable/manual/performance-tips/}
\end{frame}
\begin{frame}{About global scope}
  \begin{block}{Performance tips}
    A global variable might have its value, and therefore its type, change at any point. This makes it difficult for the compiler to optimize code using global variables. Variables should be local, or passed as arguments to functions, whenever possible.

    Any code that is performance critical or being benchmarked should be inside a function.
  \end{block}
  \pause
  \begin{block}{}
    The natural unit of compilation is a function. Code in global scope does not get compiled.
  \end{block}
\end{frame}
\begin{frame}[fragile]{Type instabilities}
  Julia relies on type-inference to generate optimal code, when your code has type-instabilities your code will execute slower.
  Some of these problems can be fixed by compiler improvements over time, but if you want the fast code now they are best avoided,
  especially in hot code. Sometimes introduction a function barrier can help generating optimal code.
  \pause
  \begin{lstlisting}
  typeinstable() = rand() < 0.5 ? 1.0 : 1
  \end{lstlisting}
\end{frame}
\begin{frame}[fragile]{Using BenchmarksTools.jl}
  \begin{block}{}
    BenchmarkTools provides a variety of tools to run benchmark and compare them. See \url{https://github.com/JuliaCI/BenchmarkTools.jl}
  \end{block}
  \begin{lstlisting}
julia> using BenchmarkTools
julia> @benchmark sin(1)
julia> @benchmark sum(rand(1000))
# Interpolate inputs and globals into your benchmark to measure the right thing.
julia> @benchmark sum(\$(rand(1000)))
  \end{lstlisting}
\end{frame}
\begin{frame}[fragile]{Using the Profiler}
  You can profile a piece of code with Julia's inbuilt profiler.
  \begin{description}
      \lstitem{@profile fun()} Profile a specific function
      \lstitem{Profile.clear()} Clear the recorded profile
      \lstitem{Profile.print()} Print the profile
      \lstitem{Profile.print(C=true)} Print the profile including stacktraces reaching into C.
  \end{description}
  \begin{block}{ProfileView.jl}
    The textual output of the profiler can be hard understand \lstinline{ProfileView.jl} gives a graphical representation.
    \begin{lstlisting}
    using ProfileView
    ProfileView.view()
    \end{lstlisting}
  \end{block}
\end{frame}
\begin{frame}[fragile]{Using the memory allocation tracker}
  To track memory allocations you have to start Julia with the memory allocation tracker enabled.
  \begin{lstlisting}
  # Track only allocation in user code
  julia --track-allocation=user
  # Track allocation in all code (includeing the Julia base)
  julia --track-allocation=all
  \end{lstlisting}
  After quiting Julia *.mem files are created that contain cumulative amounts of allocated memory.
  \begin{block}{Getting useful data}
    Since we have to start Julia with track allocations enable we will gather a lot of noisy data.
    To cut down the noise run your code in a session once and then use \lstinline{Profile.clear\_malloc\_data()}
    to reset the allocation counts and then run your code again only tracking revelant allocations.
  \end{block}
\end{frame}
\begin{frame}[fragile]{A simple example}
\begin{lstlisting}
function mysum(A)
  acc = 0
  for x in A
     acc += x
  end
  return acc
end
\end{lstlisting}
\end{frame}
\section{Performance analysis}
\begin{frame}[fragile]{A supposedly simple task}
  \begin{lstlisting}
  function myfun()
    s = 0.0
    N = 10000
    for i=1:N
        s+=det(randn(3,3))
    end
    s/N
  end
  \end{lstlisting}
\end{frame}

\begin{frame}{What is next?}
  \begin{description}
    \item[Next Session] Data Structures and Algorithms
    \item[Last Session] Parallel computing, threading, GPUs? Up to grabs.
  \end{description}
\end{frame}
\end{document}
